%% kapa: exam=1 round=34 year=2023 period=1 form=A subject=civil_law subject_ko=민법 seq=1 number=1
%% source: Q-Net 감정평가사 기출문제 (cst00309) artlSeq=5211789 file=raw/1008/5211789/2023년도 제34회 감정평가사 1차시험 1교시.pdf
%% conversion: kordoc 4.24.51 -> tools/gen1.py (auto); review=proofread-sonnet; flags=none
\begin{question}{1}
민법의 법원(法源)에 관한 설명으로 옳은 것은? (다툼이 있으면 판례에 따름)
\begin{choices}
  \item 민법 제1조에서 민법의 법원으로 규정한 ‘민사에 관한 법률’은 민법전만을 의미한다.
  \item 민법 제1조에서 민법의 법원으로 규정한 ‘관습법’에는 사실인 관습이 포함된다.
  \item 대법원이 정한 「공탁규칙」은 민법의 법원이 될 수 없다.
  \item 헌법에 의하여 체결·공포된 국제조약은 그것이 민사에 관한 것이더라도 민법의 법원이 될 수 없다.
  \item 미등기무허가 건물의 양수인에게는 소유권에 준하는 관습법상의 물권이 인정되지 않는다.
\end{choices}
\end{question}
\answer{5}
